\subsection{Zariski 环}

命题 6。设 A 为交换 Noether 环，m 为 A 的理想。下列性质等价：

\begin{enumerate}
\def\labelenumi{(\alph{enumi})}
\item
  \(\mathfrak{m}\) 包含于 A 的 Jacobson 根。
\item
  每个有限生成的 A-模关于 m-进拓扑都是 Hausdorff 的。
\item
  对每个有限生成的 A-模 E，E 的每个子模关于 E 上的 m-进拓扑都是闭的。
\item
  A 的每个极大理想关于 m-进拓扑都是闭的。
\end{enumerate}

证明 (a) 蕴含 (b)。设 m 包含于 A 的 Jacobson 根，并令 E 为有限生成的 A-模。若 \(x \in E\) 和 \(m \in m\) 满足 \((1 - m)x = 0\)，则 x = 0，因为 1 - m 在 A 中可逆。于是由（第 2 号，命题 5），E 关于 m-进拓扑是 Hausdorff 的。

证明 (b) 蕴含 (c)。设 (b) 成立。令 E 为有限生成的 A-模，F 为 E 的子模。则 E/F 关于 m-进拓扑是 Hausdorff 的，而该拓扑正是 E 上 m-进拓扑的商拓扑；所以 F 在 E 中是闭的。

显然 (c) 蕴含 (d)。最后证明 (d) 蕴含 (a)。由 (d) 可知，对于 A 的每个极大理想 a，A/a 作为 A-模关于 m-进拓扑是 Hausdorff 的。这意味着 \(\mathfrak{m}(\mathrm{A}/\mathfrak{a}) \neq \mathrm{A}/\mathfrak{a}\)，除非 A/a 上的 m-进拓扑是最粗拓扑且 A/a 化为 0；但这是荒谬的，因为 A/a 是域。因此，m 在 A/a 中的典范像是 A/a 的一个不等于 A/a 的理想，故只能为 0；于是 \(m \subset a\)，这证明了 m 包含于 A 的 Jacobson 根。

定义 2。若拓扑环 A 是交换的、Noether 的，并且存在一个关于 A 上的拓扑的定义理想 m，满足命题 6 的等价条件，则称 A 为 Zariski 环。

Zariski 环 A 必然是 Hausdorff 的（命题 6），且其拓扑的每个定义理想都包含于 A 的 Jacobson 根。

\textbf{Zariski 环的例子}\par

\begin{enumerate}
\def\labelenumi{(\arabic{enumi})}
\item
  设 A 为交换 Noether 环，m 为 A 的理想。若 A 关于 m-进拓扑是 Hausdorff 且完备的，则根据第 2 节第 13 号引理 3，A 连同此拓扑是 Zariski 环。
\item
  Zariski 环的每个商环 A/b 都是 Zariski 环，因为它是 Noether 环，并且若 m 是 A 的定义理想，则 \(\mathfrak{m}(\mathbf{A}/\mathfrak{b}) = (\mathfrak{m} + \mathfrak{b})/\mathfrak{b}\) 包含于 A/b 的 Jacobson 根（代数，第八章，第 6 节，第 3 号，命题 7）。
\item
  设 A 为 Noether 半局部环，r 为其 Jacobson 根。则赋予 A r-进拓扑后，A 是 Zariski 环。当我们将 Noether 半局部环视为拓扑环时，所指的拓扑总是这一拓扑（另有说明时除外）。
\end{enumerate}

命题 7。设 A、A' 为两个交换环，h: A → A' 为环同态。假设 A 是 Noether 环，且 A' 是有限生成的 A-模（其结构由 h 定义）。设 m 为 A 的理想，m' = mA'。则：

\begin{enumerate}
\def\labelenumi{(\roman{enumi})}
\item
  对于 \(\mathfrak{m}'\)-进拓扑在 \(\mathbf{A}'\) 上为 Hausdorff，充要条件是 \(1 + h(\mathfrak{m})\) 中的元素都不是 \(\mathbf{A}'\) 中的零因子。
\item
  若赋予 A m-进拓扑后 A 是 Zariski 环，则赋予 A' m'-进拓扑后 A' 也是 Zariski 环。
\item
  若 \(h\) 是单射（从而将 \(\mathbf{A}\) 识别为 \(\mathbf{A}'\) 的子环），则 \(\mathfrak{m}'\)-进拓扑在 \(\mathbf{A}'\) 上诱导出在 \(\mathbf{A}\) 上的 \(\mathfrak{m}\)-进拓扑。
\end{enumerate}

回忆一下，A' 上的 m'-进滤过与 A'-这个 A-模上的 m-进滤过相同（第 2 节，第 1 号，例 3）。断言 (i) 因而是第 2 号命题 5 的特殊情形，断言 (iii) 是第 2 号定理 2 的特殊情形。最后证明 (ii)。设 A 连同 m-进拓扑是 Zariski 环，令 E' 为有限生成的 A'-模；它也是有限生成的 A-模，且 E' 上的 m-进滤过与 m'-进滤过相同；因此 E' 关于 m'-进拓扑是 Hausdorff 的。最后，A-模 A' 是 Noether 的，因而环 A' 是 Noether 的，这就完成了 A' 为 Zariski 环的证明。

